下一个定理讨论 Sobolev 空间之间嵌入的紧性.

\begin{Theorem}
\label{thm:compact-successive-sobolev-orders}
设 $\Omega$ 有界, $k$ 为非负整数. 假设 $\Omega$ 具有 $(k+1)$ 阶延拓性质. 那么\[H^{k+1}(\Omega)\overset{c}{\hookrightarrow}H^k(\Omega),\]即 $H^{k+1}(\Omega)$ 紧嵌入 $H^k(\Omega)$.
\end{Theorem}

\begingroup
\postdisplaypenalty=10000
\renewcommand{\endProof}{\par\nobreak\hfill$\blacksquare$\par\vspace{0.1cm}}
\begin{Proof}
设 $E$ 为延拓算子; 总可以选择 $E$, 使其映射得到的函数支于某个紧集 $S$. 若 $u_n$ 为 $H^{k+1}(\Omega)$ 中的有界序列, 则 $Eu_n$ 为 $H^{k+1}(\mathbb{R}^m)$ 中的有界序列, 且 $Eu_n$ 的支集包含于 $S$. 将证明 $Eu_n$ 有在 $H^k(\mathbb{R}^m)$ 中收敛的子序列. 用 $M_R$ 表示乘以球 $|\bm{\xi}|<R$ 的特征函数的算子. 对 $v\in H^{k+1}(\mathbb{R}^m)$, 令 $v^R=\mathcal{F}^{-1}[M_R\mathcal{F}[v]]$. 那么容易证明存在常数 $C$, 使
\begin{equation}\label{eq:sobolev-fourier-cutoff-high-frequency-bound}
\|u-v^R\|_{k,2}\leq\dfrac{C}{R}\|u\|_{k+1,2}.
\end{equation}
此外, 对任意固定的 $R$ 及任意正整数 $l$, 有 $u^R\in C_b^l(\mathbb{R}^m)$, 且存在常数 $C(R,l)$, 使
\begin{equation}\label{eq:sobolev-fourier-cutoff-smooth-norm-bound}
\|u^R\|_{l,\infty}\leq C(R,l)\|u\|_{k+1,2}.
\end{equation}
因此, 若 $Eu_n$ 在 $H^{k+1}(\mathbb{R}^m)$ 中有界, $D$ 为包含 $S$ 的任意有界开集, 则 $(Eu_n)^R$ 有一个在 $C_b^k(\overline D)$ 中收敛的子序列 (Arzela-Ascoli 定理), 从而在 $H^k(D)$ 中收敛. 选取序列 $R_n\to\infty$ 并应用对角线论证, 可找到子序列 $n_j$, 使 $(Eu_{n_j})^{R_n}$ 对每个 $n$ 都在 $H^k(D)$ 中收敛. 结合 \eqref{eq:sobolev-fourier-cutoff-high-frequency-bound}, 这意味着 $Eu_{n_j}$ 在 $H^k(D)$ 中收敛, 从而在 $H^k(\mathbb{R}^m)$ 中收敛, 因为每个 $Eu_{n_j}$ 的支集都包含于 $D$.
\end{Proof}
\endgroup

\PDEIndexMark{idx:ehrling-equivalent-norm-introduction}作为 Theorem~\ref{thm:compact-successive-sobolev-orders} 的推论, 将建立 Ehrling 引理, 它使我们能够在 $H^k(\Omega)$ 上引入等价范数. 先给出一个抽象版本.

\begin{Theorem}{Ehrling 引理}
\label{thm:ehrling-compact-intermediate-norm}
设 $X,Y$ 和 $Z$ 为 Banach 空间. 假设 $X$ 紧嵌入 $Y$, $Y$ 连续嵌入 $Z$,\[X\overset{c}{\hookrightarrow}Y\hookrightarrow Z.\]那么对每个 $\varepsilon>0$, 存在常数 $c(\varepsilon)$, 使对每个 $x\in X$ 有
\begin{equation}\label{eq:ehrling-compact-intermediate-norm-bound}
\|x\|_Y\leq\varepsilon\|x\|_X+c(\varepsilon)\|x\|_Z
\end{equation}

\end{Theorem}

\begingroup
\postdisplaypenalty=10000
\renewcommand{\endProof}{\par\nobreak\hfill$\blacksquare$\par\vspace{0.1cm}}
\begin{Proof}
假设断言对某个 $\varepsilon_0>0$ 不成立. 那么 $X$ 中存在序列 $x_n$, 使 $\|x_n\|_X=1$, 且
\begin{equation}\label{eq:ehrling-normalized-contradiction-sequence}
\|x_n\|_Y>\varepsilon_0+n\|x_n\|_Z.
\end{equation}
由于从 $X$ 到 $Y$ 的嵌入连续, $x_n$ 在 $Y$ 中有界, 而 \eqref{eq:ehrling-normalized-contradiction-sequence} 意味着 $x_n$ 必须在 $Z$ 中收敛到 $0$. 取子序列后, 可以假设 $x_n$ 在 $Y$ 中收敛; 此时极限必须为 $0$. 但这与 \eqref{eq:ehrling-normalized-contradiction-sequence} 矛盾.
\end{Proof}
\endgroup
